\section{Discussion: readings of the board, limits, and competing explanations}\label{sec:discussion}

\subsection{Where the board and the matrix disagree}\label{sec:board-mismatch}

The drawing on the board has three lines and the matrix records two edges
\cite{whiteboard1,whiteboard2}.
Edge by edge: the pair $n_1,n_2$ is drawn (with an arrowhead at $n_2$) and has
$A_{12}=A_{21}=1$; the pair $n_2,n_3$ is drawn and has $A_{23}=A_{32}=1$; the pair
$n_1,n_3$ is drawn but has $A_{13}=A_{31}=0$ (\Cref{fig:graphs}).
The matrix is symmetric, so it gives no support to a direction on the line from $n_1$ to
$n_2$.
Following \Cref{conv:example}, this note takes the matrix as the example, because the board
says that the graph is given through its adjacency matrix \cite{whiteboard1}.
If the drawing had been meant, the adjacency matrix would be $\MDrawT$, which differs from the
board matrix exactly at the entries $(1,3)$ and $(3,1)$.
This note does not decide which of the two the lecturer intended.
The choice does not affect the answers to Q1 to Q3, which are theorems for every simple
undirected graph; it changes only the numbers in the examples.
Both graphs are connected, so the test of \Cref{thm:test} answers ``connected'' for both.

\subsection{The mark above \texorpdfstring{$A$}{A}}\label{sec:board-mark}

In both photographs the mark above $A$ in question Q1 resembles a small $n$ or a pair of
quotation marks, not clearly a $k$ \cite{whiteboard1,whiteboard2}.
The board's own letter for the number of vertices is the capital $N$, so a lower-case $n$ as an
exponent would be unusual, but it cannot be ruled out.
\Cref{conv:exponent} reads the mark as $k$ because the assignment text asks about $A^k$.
Since \Cref{thm:walkcount} holds for every positive integer exponent, it answers the board
question whichever letter was written, including an exponent equal to $N$.

\subsection{Limits of the results}\label{sec:limits}

\begin{itemize}
\item \emph{Scope.} The theorems are proved for finite simple undirected graphs.
  For directed graphs and multigraphs, the walk-counting statement is made in other sources
  \cite[Theorem~4.5.10]{grinberg2025graph}\cite[Corollary~9.3.3]{lehman_mcs_adjacency}, and
  this note makes no claim about the connectivity test there.
\item \emph{What the entries do not say.} The entries count walks; they say nothing about the
  number of paths, which needs other methods and is not treated here (\Cref{prop:warning}).
\item \emph{What the test does not say.} It answers ``connected'' or ``not connected''; it
  does not list a path between two vertices.
  \Cref{cor:components} gives the components and \Cref{rem:distance} the distances.
\item \emph{Finite checks.} The computer checks of \Cref{tab:exhaustive} cover graphs with at
  most six vertices; they support the statements but the proofs are the arguments of
  \Cref{sec:test} and Appendix~\ref{app:derivations}.
\item \emph{Sources.} The textbook was searched in full text for the adjacency matrix and read
  page by page only where listed in Appendix~\ref{app:provenance}, so the statement that it
  contains neither result is a statement about the search.
  Several sources were read only in part, as Appendix~\ref{app:provenance} records.
\end{itemize}

\subsection{Competing explanations}\label{sec:competing}

\emph{Reading the entry as a count of paths.}
Theorem~1 of one source says that $(A^l)_{ij}$ is the number of ``different paths'' of length
$l$ \cite[Theorem~1]{paramadevan2021adjacency}.
As written it is wrong for any definition of path that forbids repeated vertices or repeated
edges, by \Cref{prop:warning}: $(A^2)_{ii}$ is positive although no such path exists.
The same paper states, in its property~7, that the entry equals the number of walks
\cite[property~7]{paramadevan2021adjacency}.
This note therefore does not use that Theorem~1 as support for \Cref{thm:walkcount}; the four
sources cited there state walks.

\emph{Explaining the test by paths.}
The proof of the test in the same paper applies its Theorem~1 to $A+\dots+A^{N-1}$ and so
counts ``paths'' where walks are meant, and it passes from a nonzero entry to the existence
of a path without the walk-to-path step \cite[proof of Theorem~2]{paramadevan2021adjacency}.
The proof sketch of the $(I+A)^{N-1}$ form likewise says that the power ``aggregates all
paths'' \cite[Section~1, Theorem~2, proof sketch]{bai2025adjacency}.
Both \emph{statements} of the test are correct, as \Cref{thm:test} shows; the explanations are
incomplete, which is why the proof is written out in \Cref{sec:test}.

\emph{A different method.}
Breadth-first search gives the same answer (\Cref{rem:bfs}); eigenvalue methods are outside the
scope of this note and no claim is made about them.

\subsection{Cost}\label{sec:cost}

This remark is an estimate that counts scalar multiplications with the school method, which
uses $N^3$ of them for one product of two $N\times N$ matrices; faster methods were not examined.
Reaching $A^{N-1}$ takes $N-2$ products, which is $98{,}000{,}000$ multiplications for $N=100$.
Because \Cref{prop:iplusa} allows any exponent $m\ge N-1$, squaring $I+A$ repeatedly until the
exponent is at least $N-1$ takes $7$ squarings, or $7{,}000{,}000$ multiplications, for
$N=100$.
Breadth-first search with adjacency lists takes time proportional to the number of vertices
plus the number of edges \cite[slides~9 and 10]{schost2026bfs}; when the input is the matrix
$A$ itself it reads each of the $N^2=10{,}000$ entries once \cite[slides~3 and 6]{schost2026bfs}.
The entries of the powers also grow quickly: for the complete graph on $10$ vertices the
largest entry of $A^9$ is $38{,}742{,}049$.
The matrix test is useful because it is a clean statement about powers of $A$, not because it
is fast.

\subsection{What this note does not settle}\label{sec:open}

The board shows the question ``How to know if graph is connected or not'' but not its answer
\cite{whiteboard1,whiteboard2}, so this note makes no claim about whether the lecturer's
intended answer was the sum test, the $(I+A)^{N-1}$ form, or a search.
